\section{Screening of the phonon-assisted amplitude}\label{app:D.36}

The screened phonon-assisted amplitude carries exactly one factor $1/\varepsilon_{\mathrm{mac}}(\vect{q},\omega)$, and the screening of the final pair adds a term, the pair channel, screened at the pair's own momentum and energy. We derive the factor $1/\varepsilon_{\mathrm{mac}}$ and the pair channel, and the screened loss function that \cref{sec:13.4} uses, from second-order time-dependent mean-field theory; the band-sum theorem follows as a by-product. The level is the random-phase approximation, also called time-dependent Hartree, with an exchange--correlation kernel added where stated. The electron--hole ladder of \cref{sec:4.9}, made of interaction lines inside a loop, is left out.

\subsection*{Second-order mean-field response}

The Kohn--Sham electrons move in the one-electron Hamiltonian $h$, with eigenstates $|a\rangle$, energies $\varepsilon_a$ and occupations $f_a \in \{0, 1\}$; each label $a$ stands for a band and a crystal momentum. Their density matrix $\varrho(t)$ obeys
\begin{equation*}
  \ii\,\partial_t\varrho = \bigl[h + U(t),\,\varrho\bigr] , \qquad U(t) = U_{\mathrm{ext}}(t) + f_{\mathrm{Hxc}}\,\delta\nd(t) ,
\end{equation*}
where $U_{\mathrm{ext}}$ is the applied potential and the second term the Hartree and exchange--correlation potential of the induced density $\delta\nd$, linear in $\delta\nd$ at this level. Two fields are applied. The dark matter acts at $(\vect{q},\omega)$ through the one-body operator $\mathsf{x}$, and the phonon through $\mathsf{y}$ at frequency $\omega_2 = \mp\omega_{\vect{Q}\nu}$, the upper sign for emission; the matrix element $N^{-1/2}(n_{\vect{Q}\nu} + \frac12 \pm \frac12)^{1/2}$ of the phonon operator is taken out, as in \cref{eq:13.2-M}. Each field is first dressed by its own induced density. For the phonon the dressing gives the vertex of density-functional perturbation theory, $\mathsf{y}_{ab} = g_{ab\nu}$ (\cref{eq:6.2-screen}), evaluated statically with an error of order $(\omega_{\vect{Q}\nu}/\Eg)^2$ (\cref{app:D.27}). For the dark matter the dressing gives the screened density vertex $\mathsf{x}_{ab} = \rho^{\mathrm{scr}}_{ab}$ of \cref{eq:13.4-rhoscr}, the bare vertex of the total potential at $\vect{q} + \vect{G}$. With $\varrho = \varrho^{(0)} + \varrho^{(1)} + \varrho^{(2)} + \cdots$ and $\varrho^{(0)}_{ab} = f_a\delta_{ab}$, the first-order response to a field $F$ of frequency $\omega_i$ is
\begin{equation*}
  (\omega_i - \varepsilon_a + \varepsilon_b)\,\varrho^{F}_{ab} = F_{ab}\,(f_b - f_a) , \qquad \varrho^{F}_{ab} = \frac{F_{ab}\,(f_b - f_a)}{\omega_i - \varepsilon_a + \varepsilon_b + \ii0^+} ,
\end{equation*}
the retarded solution. The part of $\varrho^{(2)}$ bilinear in the two fields oscillates at $\omega_f = \omega + \omega_2 = \omega \mp \omega_{\vect{Q}\nu}$ and obeys
\begin{equation}
  (\omega_f - \varepsilon_a + \varepsilon_b)\,\varrho^{(2)}_{ab} = \mathfrak{s}^{(2)}_{ab} + U^{(2)}_{ab}\,(f_b - f_a) ,
  \qquad
  \mathfrak{s}^{(2)}_{ab} = \sum_m\bigl[\mathsf{x}_{am}\varrho^{\mathsf{y}}_{mb} - \varrho^{\mathsf{y}}_{am}\mathsf{x}_{mb} + \mathsf{y}_{am}\varrho^{\mathsf{x}}_{mb} - \varrho^{\mathsf{x}}_{am}\mathsf{y}_{mb}\bigr] ,
  \label{eq:D.36-second}
\end{equation}
where $\mathfrak{s}^{(2)}$ is the source that each field builds on the first-order response to the other, and $U^{(2)}$ the potential of the second-order density. The sum over $m$ runs over every state, and $\mathfrak{s}^{(2)}_{ab}$ is nonzero for every kind of pair $(a, b)$: empty and filled, filled and filled, empty and empty.

\emph{The band-sum theorem again.} Take $a = c$ empty and $b = v$ filled. For $m$ empty, $\varrho^{\mathsf{y}}_{mv} = \mathsf{y}_{mv}/(\omega_2 + \varepsilon_v - \varepsilon_m)$ and $\varrho^{\mathsf{y}}_{cm} = 0$, and likewise for $\mathsf{x}$; the source collects $\mathsf{x}_{cm}\mathsf{y}_{mv}/(\omega_2 + \varepsilon_v - \varepsilon_m) + \mathsf{y}_{cm}\mathsf{x}_{mv}/(\omega + \varepsilon_v - \varepsilon_m)$, the electron paths. For $m$ filled, $\varrho^{\mathsf{y}}_{mv} = 0$ and $\varrho^{\mathsf{y}}_{cm} = \mathsf{y}_{cm}/(\omega_2 + \varepsilon_m - \varepsilon_c)$, and the source collects $-\mathsf{y}_{cm}\mathsf{x}_{mv}/(\omega_2 + \varepsilon_m - \varepsilon_c) - \mathsf{x}_{cm}\mathsf{y}_{mv}/(\omega + \varepsilon_m - \varepsilon_c)$, the hole paths with their sign. On the shell $\omega_f = \varepsilon_c - \varepsilon_v$ the hole denominators are $\omega_2 + \varepsilon_m - \varepsilon_c = -(\omega + \varepsilon_v - \varepsilon_m)$ and $\omega + \varepsilon_m - \varepsilon_c = -(\omega_2 + \varepsilon_v - \varepsilon_m)$, and
\begin{equation*}
  \mathfrak{s}^{(2)}_{cv} = \Bigl\langle c\Bigl|\,\mathsf{y}\,\frac{1}{\varepsilon_v + \omega - h}\,\mathsf{x} + \mathsf{x}\,\frac{1}{\varepsilon_v + \omega_2 - h}\,\mathsf{y}\,\Bigr|v\Bigr\rangle ,
\end{equation*}
which is \cref{eq:D.35-theorem}: the source of the pair is the four-path amplitude with the unrestricted band sum, here evaluated with the screened vertex $\mathsf{x} = \rho^{\mathrm{scr}}$. The terms $m = v$ and $m = c$ are included, with the denominators $\omega$ and $-\omega$ of \cref{app:D.35}. The density-matrix derivation needs no Fock space: the filled-band terms enter through the occupation differences.

\subsection*{The pair channel}

The second-order density at the wavevector $\vect{q} + \vect{Q} + \vect{G}$ is the trace of $\varrho^{(2)}$ with $\ee^{-\ii(\vect{q}+\vect{Q}+\vect{G})\cdot\vect{r}}$, two spins included. With $b$ at $\vect{k}$ and $a$ at $\vect{k}+\vect{q}+\vect{Q}$, the matrix element of the plane wave is $\rho_{ab}(\vect{k},\vect{q}+\vect{Q}+\vect{G})^*$. \Cref{eq:D.36-second} then gives two terms:
\begin{equation*}
  \delta\nd^{(2)}_{\mathrm{tot},\vect{G}} = \delta\nd^{(2)}_{\vect{G}} + \sum_{\vect{G}'}R^0_{\vect{G}\vect{G}'}(\vect{q}+\vect{Q},\omega_f)\,U^{(2)}_{\vect{G}'} ,
\end{equation*}
where $R^0$ is the bubble \cref{eq:4.6-R0}, since $U^{(2)}_{ab} = \sum_{\vect{G}'}U^{(2)}_{\vect{G}'}\rho_{ab}(\vect{k},\vect{q}+\vect{Q}+\vect{G}')$, and $\delta\nd^{(2)}_{\vect{G}}(\vect{q}+\vect{Q})$ of \cref{eq:13.4-dn2} is the density that the source alone creates in the Kohn--Sham electrons, summed over all pairs of states. The potential $U^{(2)}$ closes the loop, $U^{(2)} = f_{\mathrm{Hxc}}\,\delta\nd^{(2)}_{\mathrm{tot}}$, a matrix product at $(\vect{q}+\vect{Q},\omega_f)$. Eliminating $\delta\nd^{(2)}_{\mathrm{tot}}$ gives
\begin{equation*}
  U^{(2)} = f_{\mathrm{Hxc}}\bigl(1 - R^0f_{\mathrm{Hxc}}\bigr)^{-1}\delta\nd^{(2)} = \bigl(1 - f_{\mathrm{Hxc}}R^0\bigr)^{-1}f_{\mathrm{Hxc}}\,\delta\nd^{(2)} = \tilde\varepsilon^{-1}f_{\mathrm{Hxc}}\,\delta\nd^{(2)} ,
\end{equation*}
with $\tilde\varepsilon^{-1} = (1 - f_{\mathrm{Hxc}}R^0)^{-1} = 1 + f_{\mathrm{Hxc}}R$ the inverse dielectric matrix of a test electron (\cref{sec:8.7}), retarded, at $(\vect{q}+\vect{Q},\omega_f)$. In the random-phase approximation $\tilde\varepsilon^{-1}f_{\mathrm{Hxc}} = \varepsilon^{-1}v_{\mathrm{C}}$ is the retarded screened interaction $W^{\mathrm{scr,R}}(\vect{q}+\vect{Q},\omega_f)$ of \cref{eq:4.8-W}.

For the final pair, $a = c$ at $\vect{k}+\vect{q}+\vect{Q}$ and $b = v$ at $\vect{k}$, \cref{eq:D.36-second} reads $(\omega_f - \varepsilon_c + \varepsilon_v)\varrho^{(2)}_{cv} = \mathfrak{s}^{(2)}_{cv} + U^{(2)}_{cv}$. This equation for $\varrho^{(2)}_{cv}$ has the form of a first-order coherence driven at $\omega_f$ by the one-body potential $\mathfrak{s}^{(2)} + U^{(2)}$, and the numerator of its solution is the amplitude for creating the pair, exactly as $F_{cv}$ is for a single field. The bare pairs are the asymptotic states of the mean-field evolution, and the potential that creates them is the retarded potential (\cref{sec:8.7}). The amplitude for creating the pair is therefore
\begin{equation}
  \mathfrak{s}^{(2)}_{cv} + U^{(2)}_{cv} = \mathfrak{s}^{(2)}_{cv} + \sum_{\vect{G}\vect{G}'}\rho_{cv}(\vect{k},\vect{q}+\vect{Q}+\vect{G})\,\bigl[\tilde\varepsilon^{-1}f_{\mathrm{Hxc}}\bigr]_{\vect{G}\vect{G}'}(\vect{q}+\vect{Q},\omega_f)\,\delta\nd^{(2)}_{\vect{G}'}(\vect{q}+\vect{Q}) ,
  \label{eq:D.36-amp}
\end{equation}
with $\mathfrak{s}^{(2)}_{cv}$ the four-path amplitude of the screened vertex. The pair channel is the second term.

\subsection*{The macroscopic factor}

Every term of \cref{eq:D.36-amp} is linear in $\mathsf{x} = \rho^{\mathrm{scr}}$, since $\mathfrak{s}^{(2)}$ is linear in $\mathsf{x}$. The external screening therefore factors out of the whole amplitude once it factors out of $\rho^{\mathrm{scr}}$. Split the kernel into its macroscopic Coulomb term, $v_0 = v_{\mathrm{C}}(q)$ at $\vect{G} = \vect{G}' = 0$, and the remainder $\bar f = f_{\mathrm{Hxc}} - v_0 = \bar v + f_{\mathrm{xc}}$ with $\bar v$ the Coulomb kernel without its head (\cref{sec:4.6}). Let $\bar R$ be the response to $\bar f$ alone, $\bar R = R^0 + R^0\bar f\bar R$, which is the response of \cref{sec:4.4} with the macroscopic Coulomb term removed. Then $R = \bar R + \bar Rv_0R$, so that $R_{\vect{G}0} = \bar R_{\vect{G}0}(1 + v_0R_{00})$, and
\begin{equation*}
  \bigl(\tilde\varepsilon^{-1}\bigr)_{\vect{G}0} = \delta_{\vect{G}0} + v_0\delta_{\vect{G}0}R_{00} + \bigl(\bar fR\bigr)_{\vect{G}0} = \bigl(1 + v_0R_{00}\bigr)\bigl[\delta_{\vect{G}0} + \bigl(\bar f\bar R\bigr)_{\vect{G}0}\bigr] .
\end{equation*}
The prefactor is $(\varepsilon^{-1})_{00} = 1/\varepsilon_{\mathrm{mac}}$ with $\varepsilon_{\mathrm{mac}} = 1 - v_0\bar R_{00}$ (\cref{eq:4.4-epsmac}). In the random-phase approximation $\bar f\bar R$ has the elements $v_{\mathrm{C}}(|\vect{q}+\vect{G}|)\bar R_{\vect{G}0}$ for $\vect{G} \ne 0$ and none at $\vect{G} = 0$, which gives \cref{eq:13.4-rhoLF}: $\rho^{\mathrm{scr}} = \rho^{\mathrm{LF}}/\varepsilon_{\mathrm{mac}}(\vect{q},\omega)$ exactly. The whole amplitude is $\tilde M^\pm/\varepsilon_{\mathrm{mac}}(\vect{q},\omega)$, with $\tilde M^\pm$ the amplitude \cref{eq:D.36-amp} built with $\mathsf{x} = \rho^{\mathrm{LF}}$, as in \cref{eq:13.4-Mscr}. With an exchange--correlation kernel, $\bar f$ adds the elements of $f_{\mathrm{xc}}$, including its head, and $\rho^{\mathrm{LF}}$ gains the corresponding terms; the factorization is unchanged.

\subsection*{The rate}

The transition rate into the pair is the golden rule with the amplitude \cref{eq:D.36-amp}, summed over final states as in \cref{sec:13.3}, and it gives \cref{eq:13.3-LPA} with $M^\pm \to \tilde M^\pm/\varepsilon_{\mathrm{mac}}$. The optical theorem of the mean-field evolution states that the sum over bare pairs is the whole absorptive part of the one-phonon response, with nothing counted twice. The one-phonon response at $(\vect{q},\omega)$ with its external chains removed is a sum of two kinds of diagram (\cref{sec:8.4-lines}): those with both dark-matter vertices on one loop, whose cut at the pair poles gives $\sum_f|\mathfrak{s}^{(2)}_f|^2$, and those with a pair segment, two second-order densities \cref{eq:13.4-dn2} joined by $\tilde\varepsilon^{-1}f_{\mathrm{Hxc}}$. The absorptive part of the second kind has three pieces. The cut of either density at a pair pole gives $\rho_f^*\mathfrak{s}^{(2)}_f$; the absorptive part of the screened interaction is $W^{\mathrm{scr,R}\dagger}\,\bigl[\tfrac{1}{2\ii}(R^0 - R^{0\dagger})\bigr]\,W^{\mathrm{scr,R}}$, by \cref{eq:8.7-factor} multiplied by $f_{\mathrm{Hxc}}$ on both sides, and the cut of $R^0$ is again a sum over bare pairs. The cross pieces and the square of the pair term complete the square $|\mathfrak{s}^{(2)}_f + U^{(2)}_f|^2$ pair by pair when the interaction is retarded. The advanced interaction would give the same sum with a different share per pair, as in \cref{sec:8.7}.

Inside the window the external chains are not absorptive at this order, $\varepsilon_{\mathrm{mac}}(\vect{q},\omega)$ is real up to terms of order $g^2$ (\cref{sec:13.4}), and \cref{eq:13.4-Lscr} follows with an error $O(g^4)$. Outside the window $\varepsilon_2^0 \ne 0$, and \cref{eq:13.4-Lscr} fails (\cref{sec:13.6}).

\subsection*{Two solvable limits}

\emph{A source proportional to a potential.} Suppose the source acts on every pair at $\vect{q} + \vect{Q}$ as an applied potential $\Phi$ would, $\mathfrak{s}^{(2)}_{ab} = \Phi\,\rho_{ab}(\vect{k},\vect{q}+\vect{Q})(f_b - f_a)$ with $\vect{G} = 0$ only. Then $\delta\nd^{(2)}_{\vect{G}} = R^0_{\vect{G}0}\Phi$, and with $1 + \tilde\varepsilon^{-1}f_{\mathrm{Hxc}}R^0 = \tilde\varepsilon^{-1}$,
\begin{equation*}
  \mathfrak{s}^{(2)}_{cv} + U^{(2)}_{cv} = \Phi\sum_{\vect{G}}\rho_{cv}(\vect{k},\vect{q}+\vect{Q}+\vect{G})\,\bigl(\tilde\varepsilon^{-1}\bigr)_{\vect{G}0}(\vect{q}+\vect{Q},\omega_f) ,
\end{equation*}
the source screened like an external potential at the pair's momentum and energy. Without local fields the screened source is $\mathfrak{s}^{(2)}_{cv}/\varepsilon(\vect{q}+\vect{Q},\omega_f)$.

\emph{A rigid translation.} At $\vect{Q} = 0$ a uniform displacement $u(t)$ of all nuclei couples through $\mathsf{y} = u[\mathcal{A}, h]$ with $\mathcal{A} = -\ii\ell\,\hat{\vect{e}}\cdot\op{\vect{p}}$, the generator of translations of the electrons (\cref{eq:6.2-asr}), and $[\ee^{\ii\vect{q}\cdot\vect{r}}, \mathcal{A}] = c_{\mathcal{A}}\ee^{\ii\vect{q}\cdot\vect{r}}$ with $c_{\mathcal{A}} = \ii\ell\,\hat{\vect{e}}\cdot\vect{q}$. In the frame that moves with the nuclei, $\ee^{u\mathcal{A}}$ removes $\mathsf{y}$ to first order in $u$, multiplies the dark-matter potential by $1 + c_{\mathcal{A}}u$, and adds $-\ii\dot u\mathcal{A}$. Since $\mathcal{A}$ commutes with the Coulomb interaction and with the Hartree potential, the term $c_{\mathcal{A}}u\,\phi\,\ee^{\ii\vect{q}\cdot\vect{r}}$ at frequency $\omega_f$ is a perturbation of the undisplaced crystal, screened by $\varepsilon(\vect{q},\omega_f)$; the term in $\dot u$ is of order $\omega_2$. For pairs at $(\vect{q},\omega_f)$, without local fields, the amplitude is $c_{\mathcal{A}}\rho_{cv}(\vect{k},\vect{q})/\varepsilon(\vect{q},\omega_f) + O(\omega_2)$. For the four-path term alone, with $\mathsf{y}_{ab} = u(\varepsilon_b - \varepsilon_a)\mathcal{A}_{ab}$, write $\varepsilon_p - \varepsilon_c = (\varepsilon_p - \varepsilon_v - \omega) - \omega_2$ in the first term of \cref{eq:D.35-theorem} and $\varepsilon_v - \varepsilon_p = (\varepsilon_v - \varepsilon_p + \omega_2) - \omega_2$ in the second; the parts without $\omega_2$ cancel their denominators and sum to $u[\mathsf{x}, \mathcal{A}]_{cv} = c_{\mathcal{A}}u\,\mathsf{x}_{cv}$, so $\mathfrak{s}^{(2)}_{cv} = c_{\mathcal{A}}u\,\rho_{cv}/\varepsilon(\vect{q},\omega) + O(\omega_2)$, screened at the energy deposited. The Kohn--Sham density of the source follows from the same frame argument without self-consistency: it is the response $c_{\mathcal{A}}u\,\phi'R^0(\vect{q},\omega_f)$ to the phase factor, minus the first-order density $c_{\mathcal{A}}u\,\phi'R^0(\vect{q},\omega)$ carried along by the translation, with $\phi' = \phi/\varepsilon(\vect{q},\omega)$ the screened potential. The pair term of \cref{eq:D.36-amp} is then
\begin{equation*}
  \rho_{cv}\,\frac{v_{\mathrm{C}}(q)}{\varepsilon(\vect{q},\omega_f)}\,\frac{c_{\mathcal{A}}u\bigl[R^0(\vect{q},\omega_f) - R^0(\vect{q},\omega)\bigr]}{\varepsilon(\vect{q},\omega)}
  = c_{\mathcal{A}}u\,\rho_{cv}\Bigl[\frac{1}{\varepsilon(\vect{q},\omega_f)} - \frac{1}{\varepsilon(\vect{q},\omega)}\Bigr] ,
\end{equation*}
using $v_{\mathrm{C}}R^0 = 1 - \varepsilon$, and the four-path term and the pair term add to the frame result.

\subsection*{The adiabatic local-density source}

An adiabatic local kernel takes $v_{\mathrm{xc}}$ as a function of the local density, with the expansion \cref{eq:13.4-kxc} to second order. With $\delta\nd = \delta\nd_{\mathsf{x}} + \delta\nd_{\mathsf{y}} + \delta\nd^{(2)}_{\mathrm{tot}}$, the first-order densities induced by the dark matter at $(\vect{q},\omega)$ and by the phonon at $(\vect{Q},\omega_2)$, the part bilinear in the two fields is $f_{\mathrm{xc}}\delta\nd^{(2)}_{\mathrm{tot}} + k_{\mathrm{xc}}\,\delta\nd_{\mathsf{x}}\delta\nd_{\mathsf{y}}$; the cross term of $\frac12(\delta\nd_{\mathsf{x}} + \delta\nd_{\mathsf{y}})^2$ has coefficient one. The term $f_{\mathrm{xc}}\delta\nd^{(2)}_{\mathrm{tot}}$ is in $U^{(2)}$ already. The term $k_{\mathrm{xc}}\,\delta\nd_{\mathsf{x}}\delta\nd_{\mathsf{y}}$ is a new potential at $(\vect{q}+\vect{Q}+\vect{G},\omega_f)$, the Fourier component of a product in real space, which the electrons screen as any applied potential:
\begin{equation*}
  U^{(2)} = \tilde\varepsilon^{-1}\bigl[f_{\mathrm{Hxc}}\,\delta\nd^{(2)} + (k_{\mathrm{xc}}\,\delta\nd_{\mathsf{x}}\delta\nd_{\mathsf{y}})\bigr] ,
\end{equation*}
which adds $\sum_{\vect{G}}\rho_{cv}(\vect{k},\vect{q}+\vect{Q}+\vect{G})\,[\tilde\varepsilon^{-1}(k_{\mathrm{xc}}\delta\nd_{\mathsf{x}}\delta\nd_{\mathsf{y}})]_{\vect{G}}$ to \cref{eq:D.36-amp}. For exchange alone, $\nd\varepsilon_{\mathrm{x}} = -C_{\mathrm{x}}\nd^{4/3}$ gives $f_{\mathrm{x}} = -\frac49C_{\mathrm{x}}\nd^{-2/3}$ and $k_{\mathrm{x}} = \frac{8}{27}C_{\mathrm{x}}\nd^{-5/3} = -\frac23 f_{\mathrm{x}}/\nd$. Against the Hartree part of the pair term, $v_{\mathrm{C}}\,\delta\nd^{(2)}$, the new term is then $\frac23|f_{\mathrm{x}}/v_{\mathrm{C}}(|\vect{q}+\vect{Q}|)|$ times $|\delta\nd_{\mathsf{x}}\delta\nd_{\mathsf{y}}|/(\nd\,|\delta\nd^{(2)}|)$, which with the gas value $f_{\mathrm{x}} = -e^2/4k_{\mathrm{F}}^2$ (\cref{sec:4.6}) is the ratio quoted in \cref{sec:13.4}. The last factor is of order one in the electron gas, whose second-order response is of order $R^2/\nd$.
